\section{Related Work}
\label{sec:related}

We organize prior work by the three axes STOA unifies; \Cref{tab:positioning} summarizes which axis each
line decides and by what means.

\paragraph{Representation and memory-OS.}
MemOS~\cite{F1_li2025memos} and MemGPT~\cite{P14_packer2023memgpt} treat memory as an operating-system
problem, exposing abstractions (MemCube, virtual context) and primitives to migrate items between
plaintext, activation, and parameter forms; M+~\cite{F2_wang2025mplus} adds a co-trained latent memory.
These provide the \emph{mechanism} to change representation and tier but decide \emph{when} by heuristic;
STOA supplies the learned scheduler they lack.

\paragraph{Content and RL-based management.}
Memory-R1~\cite{F6_yan2025memoryr1} and Memory-as-Action~\cite{F7_zhang2025memoryasaction} learn, with
RL, \emph{what} to add, update, or delete and how to edit context, and the latter handles the
trajectory-fracturing that edits cause. But they operate over a single fixed representation and never
choose a physical tier or schedule offline consolidation---orthogonal to, and composable with, STOA.

\paragraph{Systems tiering and ML-for-systems.}
LMCache~\cite{F5_liu2025lmcache} tiers KV-cache bytes to cut miss-rate; H2O~\cite{S1_zhang2023h2o} evicts
by attention score. The learnability of such decisions is established: PARROT~\cite{S2_liu2020parrot}
imitates a Belady oracle for cache replacement, Cold-RL~\cite{S4_coldrl2025} does eviction with offline
RL in production, and Decima~\cite{S5_mao2019decima} schedules clusters with a GNN policy. STOA inherits
these techniques but changes the objective from miss-rate to downstream \emph{task utility}, and the
action from a single axis to the joint representation\,$\times$\,tier\,$\times$\,timing decision.

\paragraph{Offline consolidation.}
Sleep-time compute~\cite{F3_lin2025sleeptime} and LightMem~\cite{F4_fang2025lightmem} show the value of
precomputing and consolidating memory offline, but treat it as a global mode rather than a per-item,
budget-aware timing decision---the third axis STOA learns.

\paragraph{Learned and adaptive caching.}
Two lines bear directly on our results and we position against them explicitly, because our findings
reproduce rather than extend them. \textbf{LRB}~\cite{S6_song2020lrb} learns a relaxed-Belady objective
for CDN caching from a feature vector built out of the deltas between an object's past request
times --- the same family as the think-time features we tried here. Timing as a cache-scheduler input is
therefore established practice, not a new interface insight, and our (negative) placement result should
be read as a report that this family does not transfer to KV-block \emph{placement} under our cost model,
not as a novel signal. \textbf{ARC}~\cite{S7_megiddo2003arc} adaptively balances recency against
frequency precisely because neither dominates across workloads; the recency-versus-frequency inversion we
measure between our two sources (\Cref{sec:reaction}) is that observation restated on KV blocks, and an
adaptive policy of the ARC family is the baseline a deployment should reach for rather than either pure
rule.

The rest of the learned-cache literature divides by what it learns \emph{from} and what it decides, and
the division explains why our result is not a contradiction of it.
LeCaR~\cite{S8_vietri2018lecar} and CACHEUS~\cite{S9_rodriguez2021cacheus} learn no per-item model at all:
they run experts (LRU/LFU, and in CACHEUS scan- and churn-resistant variants) and update expert weights
from regret observed on a ghost list. That is the same evidence channel ARC uses, and
\Cref{sec:reaction} measures how thin it is here --- 3.5--7.3\% of misses carry any signal --- so the
family inherits the limitation rather than escaping it. We ran both to check, rather than leaving it as a
prediction, and both land in the same few-point band as ARC. LeCaR is nonetheless the best single policy
at the tightest tier on both traces, which is the one place the regret formulation visibly pays.
Hawkeye~\cite{S10_jain2016hawkeye} and Mockingjay~\cite{S11_shah2022mockingjay} are the closest
antecedents to what we attempted: both supervise a predictor with Belady replayed over history, Hawkeye
as a binary cache-friendly/averse label and Mockingjay as a multiclass reuse-distance estimate. They
operate at CPU last-level-cache granularity, where a program counter identifies a reuse pattern; a KV
block hash carries no analogous instruction identity, which is the structural reason their signal does
not simply port. Baleen~\cite{S12_wong2024baleen} is the one that decides admission and prefetching
rather than eviction, and the one that optimizes a downstream system cost (disk-head time) instead of
miss rate --- the closest precedent for the multi-objective reward of \Cref{sec:formulation}, and the
direction we think is right even though we do not evaluate it here.

Of these we evaluate LRB, LeCaR and CACHEUS (\Cref{sec:reaction}) --- and the prediction above is what we
measured: both expert policies land in the same few-point band as ARC rather than above it. Hawkeye,
Mockingjay and Baleen are positioned but not run.

In short, each line optimizes one axis and freezes the others, and none decides representation, tier,
\emph{and} timing jointly under a shared budget (\Cref{tab:positioning}). Our contribution is not a better
policy but a measurement: on production KV traces, ranking quality improves substantially while the
placement benefit it buys does not, and we report the protocol needed to see that.
